\section*{الفصل \ref{ch:b2:wittig} --- تفاعل فيتيغ وتفاعلات أخرى مكوِّنة للرابطة C=C}

\begin{solution}{exo:b2:wittig:1}
\ce{Ph3P + CH3CH2Br -> [Ph3PCH2CH3]+ + Br-} (استبدال ثنائي الجزيئية، بالتسخين في مذيب)، ثم
البيوتيل ليثيوم في إيثر جاف أو في THF تحت غاز خامل: \ce{[Ph3PCH2CH3]+ + BuLi -> Ph3P=CHCH3 + BuH + Li+}.
\end{solution}

\begin{solution}{exo:b2:wittig:2}
يعطي الإيليد \ce{Ph3P=CH2} والبنزالدهيد الفينيل إيثين (الستيرين)، \ce{C6H5CH=CH2}، 
وأكسيد ثلاثي فينيل الفوسفين.
\end{solution}

\begin{solution}{exo:b2:wittig:3}
\ce{Ph3P=CHCOOEt} و\ce{Ph3P=CHCOCH3}: مثبتان. \ce{Ph3P=CHCH3}: غير مثبت، وهو الذي
يحتاج إلى البيوتيل ليثيوم (أو هيدريد الصوديوم، أو أميد الصوديوم). \ce{Ph3P=CHC6H5}: مثبت بأريل، يُحضَّر عمليًا
بألكوكسيد أو بهيدروكسيد، ويعطي خلائط من الهندستين.
\end{solution}

\begin{solution}{exo:b2:wittig:4}
(\textit{E})-بنت-2-ينوات الإيثيل، \ce{CH3CH2CH=CHCOOEt}، مع فوسفات ثنائي الإيثيل الصوديومي.
\end{solution}

\begin{solution}{exo:b2:wittig:5}
فيتيغ: الإيليد \ce{Ph3P=CH2}، المحضَّر من يوديد ميثيل ثلاثي فينيل الفوسفونيوم والبيوتيل ليثيوم، يعطي مع الهكسانون الحلقي
ميثيلين هكسان حلقي وحده، والرابطة المزدوجة خارج الحلقة. أما نزع الماء من 1-ميثيل هكسان حلقي-1-ول فيمر
عبر كاتيون كربوني ثالثي ويفقد الهيدروجين الذي يعطي الألكين الأكثر استبدالًا (زايتسيف):
فيعطي أساسًا 1-ميثيل هكسين حلقي، وهو المتماكب الخطأ.
\end{solution}

\begin{solution}{exo:b2:wittig:6}
القطعان متماثلان: البنزالدهيد وإيليد البنزيليدين \ce{Ph3P=CHC6H5}. وهذا الإيليد ليس
مثبتًا إلا بأريل فيعطي خليطًا (\textit{E})/(\textit{Z}). وللحصول على (\textit{E}) النقي، استعمل تفاعل HWE
لبنزيل \omterm{def:b2:wittig:hwe}{فوسفونات} ثنائي الإيثيل مع البنزالدهيد.
\end{solution}

\begin{solution}{exo:b2:wittig:7}
مع \ce{Ph3P=CHCH3}: (\textit{Z})-بنت-2-ين، \ce{CH3CH2CH=CHCH3}، و(\textit{Z}) غالبًا. ومع
\ce{Ph3P=CHCOOEt}: (\textit{E})-بنت-2-ينوات الإيثيل، و(\textit{E}) غالبًا.
\end{solution}

\begin{solution}{exo:b2:wittig:8}
\ce{Ph3P=CH2 + C6H10O -> C7H12 + Ph3PO}. الكتل: ميثيلين هكسان حلقي \qty{96.173}{g/mol}،
وأكسيد ثلاثي فينيل الفوسفين \qty{278.291}{g/mol}؛ والاقتصاد الذري $96.173/(96.173 + 278.291) =
96.173/374.464 \approx 25.7~\%$. فثلاثة أرباع كتلة الكواشف تخرج على هيئة أكسيد الفوسفين.
% ledger: aw:C, aw:H, aw:O, aw:P
\end{solution}

\begin{solution}{exo:b2:wittig:9}
الفوسفور، وهو محب للنواة، يزيح البروميد:
\begin{equation*}
\ce{P(OEt)3 + BrCH2COOEt -> [(EtO)3PCH2COOEt]+ + Br-}.
\end{equation*}
ثم يهاجم
البروميد كربونًا من إحدى مجموعات الإيثيل في أيون الفوسفونيوم (استبدال ثنائي الجزيئية)، فيحرر
البروموإيثان ويكوّن الرابطة \ce{P=O}:
\begin{equation*}
\ce{[(EtO)3PCH2COOEt]+ + Br- -> (EtO)2P(O)CH2COOEt + EtBr}.
\end{equation*}
\end{solution}

\begin{solution}{exo:b2:wittig:10}
\ce{OHC-CH=CH-CHO + 2Ph3P=CHC6H5 -> C6H5CH=CH-CH=CH-CH=CHC6H5 + 2Ph3PO}: مكافئان من
إيليد البنزيليدين (من كلوريد بنزيل ثلاثي فينيل الفوسفونيوم وقاعدة)، واحد لكل مجموعة ألدهيد.
وتتكوّن الروابط المزدوجة الجديدة خلائط من الهندستين؛ ويُعزل المتماكب (\textit{E}) الكلي، الأكثر
استقرارًا، بالتبلور، أو يُماكَب الخليط.
\end{solution}

\begin{solution}{exo:b2:wittig:11}
الألدولي: البنزالدهيد مع البروبانون بفائض وهيدروكسيد الصوديوم (\cref{ch:b2:enolates-aldol})؛
كواشف رخيصة، والماء ناتج ثانوي، والناتج (\textit{E})، لكن التكاثف قد يجري مرتين
(ثنائي بنزال الأسيتون) إن لم يكن البروبانون بفائض. وفيتيغ: البنزالدهيد مع \omterm{def:b2:wittig:ylide}{الإيليد المثبت}
\ce{Ph3P=CHCOCH3}؛ (\textit{E})، وناتج واحد، بلا تكاثف مزدوج، لكن الإيليد مكلف 
ويجب إزالة أكسيد ثلاثي فينيل الفوسفين. والطريق الألدولي هو الأفضل هنا.
\end{solution}

\begin{solution}{exo:b2:wittig:12}
فيتيغ: البيوتانال مع \ce{Ph3P=CHCH2CH2CH3} (من 1-بروموبيوتان، ثم البيوتيل ليثيوم)، في ظروف خالية من الأملاح:
(\textit{Z}) غالبًا؛ والناتج الثانوي أكسيد ثلاثي فينيل الفوسفين، وشيء من المتماكب (\textit{E}). والألكاين:
أوكت-4-اين (من الإيثاين بألكلتين بواسطة 1-بروموبروبان وأميد الصوديوم) مهدرجًا على
حفاز بلاديوم مسموم (\cref{prop:b2:alkene-redox:syn-hydrogenation}): (\textit{Z}) وحده؛ 
والنواتج الثانوية بروميد الصوديوم والأمونيا من الألكلتين.
\end{solution}

\begin{solution}{pb:b2:wittig:1}
\textbf{1.} \ce{C23H46}، أي \ce{C_nH_{2n}} مع $n = 23$: رابطة مزدوجة واحدة (أو حلقة واحدة)، وهي هنا
\ce{C=C}.
\textbf{2.} الرابطة المزدوجة \ce{C9=C10}.
\textbf{3.} النونانال \ce{CH3(CH2)7CHO} مع الإيليد المحضَّر من 1-هالوتتراديكان؛ أو التتراديكانال
\ce{CH3(CH2)12CHO} مع الإيليد المحضَّر من 1-هالونونان.
\textbf{4.} 1-بروموتتراديكان، \ce{CH3(CH2)13Br}.
\textbf{5.} الاستبدال الثنائي الجزيئية بثلاثي فينيل الفوسفين الضخم سريع عند كربون أولي غير معاق
ولا ينافسه الحذف.
\textbf{6.} \ce{Ph3P + C14H29Br -> [Ph3PC14H29]+ + Br-}، بروميد تتراديسيل ثلاثي فينيل الفوسفونيوم.
\textbf{7.} البيوتيل ليثيوم (أو هيدريد الصوديوم، أو أميد الصوديوم)؛ فالهيدروكسيد قاعدة أضعف من اللازم
لإيليد غير مثبت، والماء الذي يجلبه كان سيبرتن الإيليد.
\textbf{8.} لا، فهو لا يحمل إلا مجموعة ألكيل: (\textit{Z}) غالبًا (\cref{prop:b2:wittig:stereo}).
\textbf{9.} حلقة رباعية \ce{P-C-C-O}: الفوسفور مرتبط بالكربون الحامل للسلسلة
\ce{C13H27}، وذلك الكربون بالكربون الحامل للسلسلة \ce{C8H17}، وهذا بالأكسجين، المرتبط
بالفوسفور.
\textbf{10.} البيوتيل ليثيوم والإيليد يتفاعلان مع الماء والأكسجين.
\textbf{11.} تصل الرابطة المزدوجة كربون الكربونيل السابق للنونانال (C9 في الناتج) بكربون
الإيليد السابق (C10)، ولا تتحرك أي رابطة أخرى (\cref{prop:b2:wittig:regiospecific}).
\textbf{12.} كاتيون كربوني عند C9، قد يفقد هيدروجينًا من C8 أو C10 وقد ينتقل: خليط من
تريكوس-8-ين وتريكوس-9-ين، كل منهما على هيئة (\textit{E}) و(\textit{Z})، و(\textit{E}) في معظمه.
\textbf{13.} \ce{Ph3P=CHC13H27 + C8H17CHO -> C8H17CH=CHC13H27 + Ph3PO}.
\textbf{14.} الهدف \ce{C23H46}: \qty{322.621}{g/mol}؛ والنونانال \ce{C9H18O}: \qty{142.242}{g/mol}؛ والملح
\ce{C32H44BrP}: \qty{539.582}{g/mol}؛ و\ce{Ph3PO}، أي \ce{C18H15OP}: \qty{278.291}{g/mol}.
\textbf{15.} يزن الإيليد \ce{C32H43P} ما قدره $539.582 - 80.912 = \qty{458.670}{g/mol}$؛ ومع النونانال،
\qty{600.912}{g/mol} من الكواشف مقابل \qty{322.621}{g/mol} من الناتج: $322.621/600.912 \approx
53.7~\%$.
\textbf{16.} بقطبيته: يمر الهيدروكربون غير القطبي عبر عمود قصير من السيليكا بمُزيح
غير قطبي بينما يبقى الأكسيد؛ أو يتبلور الأكسيد خارجًا من مذيب غير قطبي بارد.
\textbf{17.} الرابطة \ce{P=O} الشديدة القوة المتكوّنة ترجح على الرابطة \ce{P-C} المنكسرة
(\cref{prop:b2:wittig:driving}).
\textbf{18.} يحتاج إلى \omterm{def:b2:wittig:hwe}{فوسفونات} فيها مجموعة ساحبة للإلكترونات على الكربون المتفاعل، وهي غائبة
عن سلسلة ألكيلية بسيطة، ويعطي الألكين (\textit{E}).
\textbf{19.} تريكوس-9-اين، \ce{CH3(CH2)7C#C(CH2)12CH3}، مع الهيدروجين على حفاز بلاديوم
مسموم (البلاديوم على كربونات الكالسيوم مع ملح للرصاص والكينولين).
\textbf{20.} دك-1-اين، \ce{CH3(CH2)7C#CH}، يُنزع بروتونه بأميد الصوديوم، ثم 1-بروموتريديكان
\ce{CH3(CH2)12Br}: $8 + 2 + 13 = 23$ ذرة كربون، والرابطة الثلاثية بين C9 وC10.
\textbf{21.} الهدرجة الإضافية كانت ستعطي التريكوسان، بلا رابطة مزدوجة؛ والحفاز المسموم
يربط الألكاين بقوة أكبر بكثير من الألكين، الذي يغادر السطح قبل إضافة ثانية.
\textbf{22.} يتطلب كلاهما خطوتين انطلاقًا من قطع متاحة. وطريق الألكاين يعطي المتماكب (\textit{Z})
على نحو نظيف ولا يترك إلا أملاحًا؛ أما طريق فيتيغ فيعطي شيئًا من المتماكب (\textit{E}) وكتلة من
أكسيد ثلاثي فينيل الفوسفين.
\textbf{23.} 100~\%: \ce{C23H44 + H2 -> C23H46}، فكل ذرة تنتهي في الناتج.
\textbf{24.} $278.291/322.621 = \boldsymbol{\approx \qty{0.86}{g}}$ من أكسيد ثلاثي فينيل الفوسفين لكل
غرام من الفيرومون.
\end{solution}
